%---------------------------------------------
% 			  PREAMBLE
%---------------------------------------------

\documentclass[12pt]{article}

\usepackage{amsmath,amssymb,amsfonts}
\usepackage{graphicx}
\usepackage{xcolor}
\usepackage[utf8]{inputenc}
\usepackage[english]{babel}
\usepackage{amsthm}
\usepackage{chngcntr}
\counterwithin{equation}{section}
\usepackage{setspace}
\setlength{\topmargin}{-.25in} \setlength{\textheight}{8.75in}
\setlength{\oddsidemargin}{0in} \setlength{\textwidth}{6.8in}
\setstretch{1.5}
\addtolength{\jot}{4mm}  %Equation spacing

\usepackage{booktabs}
\usepackage{xparse}
\usepackage{tikz}
\usetikzlibrary{calc}

% Allows to bold math and text together
\usepackage{bm}
\newcommand\bolden[1]{{\boldmath\bfseries#1}}


% When opened in Acrobat Reader, clicking on the solution will reveal/hide the solution text
\usepackage[tikz]{ocgx2}
\newcounter{ocgSolution}
\setcounter{ocgSolution}{1}

\newcommand{\solution}[1]{%
\newline
\switchocg{\theocgSolution}{Click for Solution:}
\begin{ocg}{OCG\theocgSolution}{\theocgSolution}{0}
#1
\end{ocg}
\addtocounter{ocgSolution}{1}
}



%---------------------------------------------
% 			  HEADER
%---------------------------------------------

\usepackage{fancyhdr}
\setlength{\headheight}{15.2pt}
\pagestyle{fancy}
\lhead{ \textbf { Kyle Hutchings} }
\rhead{ \textbf { Data Science } }

%---------------------------------------------
% 			  DOCUMENT
%---------------------------------------------

\begin{document}


%---------------------------------------------
% 		Module 1
%---------------------------------------------

\section{Probability}
\subsection{Introduction}

\begin{enumerate}

\item \bolden{What is Bayes' rule?} 
\solution{
    Bayes' rule gives the probability of event A happening given event B has happened:
    $$
    P(A|B) = \frac{P(B|A)P(A)}{P(B)}
    $$
}

\item \bolden{What is the Law of Total Probability?} 
\solution{
This is a way to think about the probability of an event occuring as a weighted sum or even an expectation-value-like setup:
$$
P(A) = P(A|B_1)P(B_1) + ... + P(A|B_n)P(B_n)
$$
For example, the probability that event $B_1$ will occur is $P(B_1)$.  The probability of $A$ occuring given $B_1$ has occured is weighted by $P(B_1)$.
}

\item \bolden{What is the permutation equation?} 
\solution{
The permutation equation is used when the order matters.  
Sometimes repetition can be allowed in a permutation -- in opening a 4-number dial lock, for example, you might have 1314 where 1 shows up twice.  
Sometimes repetition is not allowed, like drawing cards from a deck.  If I'm thinking about the probability of drawing two spades, then drawing one spade on my first draw will affect my probability of drawing another.
$$
\frac{n!}{(n-k)!}
$$
}

\item \bolden{Question} 
\solution{

}

\item \bolden{Question} 
\solution{

}

\item \bolden{Question} 
\solution{

}

\item \bolden{Question} 
\solution{

}

\end{enumerate}

\newpage





\end{document} 